% Figure from MATH 590 notes, section 29. Compile with the notes preamble.
\begin{tikzpicture}[scale=1.25, >=stealth,
  mid/.style={postaction={decorate}, decoration={markings, mark=at position 0.56 with {\arrow{Stealth[length=2.6mm]}}}}]
  \fill[black!4] (0,0) rectangle (2,2);
  \draw[figB, thick, mid] (0,0) -- (2,0);
  \draw[figR, thick, mid] (2,0) -- (2,2);
  \draw[figB, thick, mid] (0,2) -- (2,2);
  \draw[figR, thick, mid] (0,0) -- (0,2);
  \node[figB, below, font=\small] at (1,0) {$a$};
  \node[figR, right, font=\small] at (2,1) {$b$};
  \node[figB, above, font=\small] at (1,2) {$a$};
  \node[figR, left, font=\small] at (0,1) {$b$};
  \fill[black] (0,0) circle (1.8pt); \fill[black] (2,0) circle (1.8pt);
  \fill[black] (2,2) circle (1.8pt); \fill[black] (0,2) circle (1.8pt);
  \node[black, below left, font=\scriptsize] at (0,0) {$x_0$};
  \node[black, below right, font=\scriptsize] at (2,0) {$x_0$};
  \node[black, above right, font=\scriptsize] at (2,2) {$x_0$};
  \node[black, above left, font=\scriptsize] at (0,2) {$x_0$};
  \node[font=\scriptsize, black!60] at (1,1) {$aba^{-1}b^{-1}$};
\end{tikzpicture}
